% GENERATED FILE. Edit canonical lessons, then run npm run generate:books.
% canonical-source-hash: fnv1a64:966002601964545e
% canonical-lessons: FR-C425-camping, FR-C425-tente, FR-C425-sac-de-couchage, FR-C425-camper, FR-C425-bivouaquer

\chapter{Au camping}
\label{ch:fr-au-camping}
\hlchaptermodality{425}

\begin{chapteropening}
\textbf{By the end of this chapter:} \emph{I can talk about camping: the campsite, the tent, the sleeping bag, camping and sleeping out in the open.}

Complete the last lesson of chapter 425: I can talk about camping: the campsite, the tent, the sleeping bag, camping and sleeping out in the open.
\end{chapteropening}

\section[le camping]{le camping — the campsite; camping}

\label{lesson:FR-C425-camping}



\begin{quote}
Before the new one: say the French for to lodge, then the French for to stay for a while.
\end{quote}

\subsection*{You'll want to know: le camping}
\textbf{le camping} — \textquotedblleft{}the campsite; camping\textquotedblright{}.

\begin{quote}
Le camping est à cinq minutes de la plage.
\end{quote}

\begin{quote}
\emph{The campsite is five minutes from the beach.}
\end{quote}

In French, \emph{le camping} is also the place.

\begin{cousinweb}
\emph{camping} was borrowed from English \emph{camping} (EtymDB).
\end{cousinweb}

\begin{grammarlens}[title={What you've built}]
One more word for telling, explaining and describing at B1.
\end{grammarlens}

\subsection*{Your turn}
\begin{itemize}
\raggedright
  \item \emph{Say it:} \emph{le camping}
  \item \emph{Say it:} \emph{le camping}, once more
  \item \emph{Say it:} a sentence of your own with \emph{le camping}
  \item \emph{Recall:} say the French for to lodge, then the French for to stay for a while, then say \emph{le camping} again
\end{itemize}

\begin{tcolorbox}[breakable,colback=teal!4,colframe=teal!35!black,title={Before you move on}]
What does le camping mean? (\textquotedblleft{}The campsite; camping\textquotedblright{}.) Say it once more.
\end{tcolorbox}
\section[la tente]{la tente — the tent}

\label{lesson:FR-C425-tente}



\begin{quote}
Before the new one: say the French for to stay for a while, then the French for the campsite.
\end{quote}

\subsection*{You'll want to know: la tente}
\textbf{la tente} — \textquotedblleft{}the tent\textquotedblright{}.

\begin{quote}
Les enfants ont monté la tente tout seuls.
\end{quote}

\begin{quote}
\emph{The children put up the tent on their own.}
\end{quote}

\emph{Monter la tente} — to pitch the tent.

\begin{cousinweb}
\emph{tente} goes back to Vulgar Latin \emph{tenta} 'something stretched', from Latin \emph{tendere} 'to stretch' (Lewis \& Short).
\end{cousinweb}

\begin{grammarlens}[title={What you've built}]
One more word for telling, explaining and describing at B1.
\end{grammarlens}

\subsection*{Your turn}
\begin{itemize}
\raggedright
  \item \emph{Say it:} \emph{la tente}
  \item \emph{Say it:} \emph{la tente}, once more
  \item \emph{Say it:} a sentence of your own with \emph{la tente}
  \item \emph{Recall:} say the French for to stay for a while, then the French for the campsite, then say \emph{la tente} again
\end{itemize}

\begin{tcolorbox}[breakable,colback=teal!4,colframe=teal!35!black,title={Before you move on}]
What does la tente mean? (\textquotedblleft{}The tent\textquotedblright{}.) Say it once more.
\end{tcolorbox}
\section[le sac de couchage]{le sac de couchage — the sleeping bag}

\label{lesson:FR-C425-sac-de-couchage}



\begin{quote}
Before the new one: say the French for the campsite, then the French for the tent.
\end{quote}

\subsection*{You'll want to know: le sac de couchage}
\textbf{le sac de couchage} — \textquotedblleft{}the sleeping bag\textquotedblright{}.

\begin{quote}
Il fait froid la nuit, alors je dors dans mon sac de couchage.
\end{quote}

\begin{quote}
\emph{It is cold at night, so I sleep in my sleeping bag.}
\end{quote}

In speech, often \emph{le duvet}.

\begin{cousinweb}
\emph{sac} is inherited from Latin \emph{saccus}, from Greek \emph{sákkos} 'bag' (EtymDB); \emph{couchage} is built on \emph{coucher} 'to lie down'.
\end{cousinweb}

\begin{grammarlens}[title={What you've built}]
One more word for telling, explaining and describing at B1.
\end{grammarlens}

\subsection*{Your turn}
\begin{itemize}
\raggedright
  \item \emph{Say it:} \emph{le sac de couchage}
  \item \emph{Say it:} \emph{le sac de couchage}, once more
  \item \emph{Say it:} a sentence of your own with \emph{le sac de couchage}
  \item \emph{Recall:} say the French for the campsite, then the French for the tent, then say \emph{le sac de couchage} again
\end{itemize}

\begin{tcolorbox}[breakable,colback=teal!4,colframe=teal!35!black,title={Before you move on}]
What does le sac de couchage mean? (\textquotedblleft{}The sleeping bag\textquotedblright{}.) Say it once more.
\end{tcolorbox}
\section[camper]{camper — to camp}

\label{lesson:FR-C425-camper}



\begin{quote}
Before the new one: say the French for the tent, then the French for the sleeping bag.
\end{quote}

\subsection*{You'll want to know: camper}
\textbf{camper} — \textquotedblleft{}to camp\textquotedblright{}.

\begin{quote}
L'été dernier, nous avons campé près d'une rivière.
\end{quote}

\begin{quote}
\emph{Last summer we camped near a river.}
\end{quote}

\emph{Le campeur} — the camper.

\begin{cousinweb}
\emph{camper} is built on \emph{camp} (EtymDB), which goes back to Latin \emph{campus} 'field, open plain'.
\end{cousinweb}

\begin{grammarlens}[title={What you've built}]
One more word for telling, explaining and describing at B1.
\end{grammarlens}

\subsection*{Your turn}
\begin{itemize}
\raggedright
  \item \emph{Say it:} \emph{camper}
  \item \emph{Say it:} \emph{camper}, once more
  \item \emph{Say it:} a sentence of your own with \emph{camper}
  \item \emph{Recall:} say the French for the tent, then the French for the sleeping bag, then say \emph{camper} again
\end{itemize}

\begin{tcolorbox}[breakable,colback=teal!4,colframe=teal!35!black,title={Before you move on}]
What does camper mean? (\textquotedblleft{}To camp\textquotedblright{}.) Say it once more.
\end{tcolorbox}
\section[bivouaquer]{bivouaquer — to bivouac, to sleep out in the open}

\label{lesson:FR-C425-bivouaquer}



\begin{quote}
Before the new one: say the French for the sleeping bag, then the French for to camp.
\end{quote}

\subsection*{You'll want to know: bivouaquer}
\textbf{bivouaquer} — \textquotedblleft{}to bivouac, to sleep out in the open\textquotedblright{}.

\begin{quote}
En montagne, ils ont bivouaqué sous les étoiles.
\end{quote}

\begin{quote}
\emph{In the mountains they slept out under the stars.}
\end{quote}

\emph{Un bivouac} is one night in the open, often in the mountains.

\begin{cousinweb}
\emph{bivouaquer} is built on \emph{bivouac}, which is usually traced to Swiss or Alsatian German \emph{Biwacht} 'extra night watch'. EtymDB lists the verb without a source.
\end{cousinweb}

\begin{grammarlens}[title={What you've built}]
One more word for telling, explaining and describing at B1.
\end{grammarlens}

\subsection*{Your turn}
\begin{itemize}
\raggedright
  \item \emph{Say it:} \emph{bivouaquer}
  \item \emph{Say it:} \emph{bivouaquer}, once more
  \item \emph{Say it:} a sentence of your own with \emph{bivouaquer}
  \item \emph{Recall:} say the French for the sleeping bag, then the French for to camp, then say \emph{bivouaquer} again
\end{itemize}

\begin{tcolorbox}[breakable,colback=teal!4,colframe=teal!35!black,title={Before you move on}]
What does bivouaquer mean? (\textquotedblleft{}To bivouac, to sleep out in the open\textquotedblright{}.) Say it once more.
\end{tcolorbox}
